%!LW recipe = latexmk (XeLaTeX)
\documentclass[xelatex,a4paper,ja=standard]{bxjsreport}


% --- 数学・記号系 ---
\usepackage{amsmath,amsthm,mathtools}
\usepackage{amssymb}  %（amsfontsを自動で含む)

\usepackage{geometry}
\geometry{margin=25truemm}
\geometry{top=15truemm}
\usepackage{expl3}
\usepackage[unicode,colorlinks=true,linkcolor=blue,urlcolor=magenta, citecolor=blue]{hyperref}
\usepackage[warnings-off={mathtools-colon,mathtools-overbracket}]{unicode-math}

\usepackage{tcolorbox}

\usepackage{fontspec}


\newtheorem{definition}{定義}%[section]
\newtheorem{theorem}[definition]{定理}%[section]
\newtheorem{lemma}[definition]{補題}%[section]
\newtheorem{prop}[definition]{命題}%[section]
\newtheorem{corollary}[definition]{系}%[section]


\title{研究内容}
\author{}
\date{\today}

\begin{document}

\maketitle

%\renewcommand{\thesection}{\arabic{section}}

\begin{definition}[特殊線形群]
    \(SL(n, \mathbb{C}) = \{ A \in M_n(\mathbb{C}) \mid \det(A) = 1 \}\) \par
    特に, \(n = 2\) のとき, \(SL(2, \mathbb{C}) = \{\begin{pmatrix} a & b \\ c & d \end{pmatrix} : a, b, c, d, \in \mathbb{C} \mid det\begin{pmatrix} a & b \\ c & d \end{pmatrix} = ad - bc = 1\}\)
\end{definition}

\begin{theorem}[\(SL(2, \mathbb{C})\) の部分群]
    \(G \subset SL(2, \mathbb{C}), |G| < \infty\) のとき, \(G\) は共役を除いて以下のいずれかである.

    • 巡回群 \(\mathbb{Z}/n\mathbb{Z}, n \in \mathbb{N}\) \par
    • バイナリー二面体群 \(BD_{4n}, n \in \mathbb{N}\) \par
    • プラトン立体に対応するバイナリー群 \(BT_{24}, BO_{48}, BI_{120}\)
    (プラトン立体 : 全ての面が合同な正多面体, \par BT : バイナリー四面体群, BO : バイナリー八面体群, BI : バイナリー二十面体群)
\end{theorem}

\begin{corollary}
    位数 \(n\) の巡回群は, \(SL(2, \mathbb{C})\) の有限部分群である.
\end{corollary}

\end{document}